\section{Introduction}

New-onset postoperative atrial fibrillation (POAF) occurs after nearly one quarter of isolated coronary artery bypass grafting (CABG) procedures and substantially elevates postoperative stroke, heart failure, intensive care stay, readmission, and long-term mortality.\cite{aranki1996predictors,villareal2004,elchami2010poaf,lapar2014poaf,dobrev2019postoperative,jawitz2020poaf,caldonazo2023meta} Because high-impact preventive interventions---including intraoperative posterior left pericardiotomy and targeted perioperative prophylaxis---must be decided before skin incision,\cite{gaudino2021palacs,chatterjee2026sts} clinical translation depends on pre-incision prediction models that deliver calibrated risk estimates both accurately and rapidly at the point of care.

Systematic appraisals under PROBAST standards demonstrate that published POAF models suffer substantial external generalization decay, falling from an apparent median ROC-AUC of 0.71 to an external median of 0.61, with over 80\% judged at high risk of bias.\cite{fields2025poaf,lopeztuesta2026poaf,pandey2023poaf,wolff2019probast,moons2025probastai} Two methodological gaps hinder trustworthy clinical benchmarking in national surgical registries. First, center-level POAF incidence spans a threefold range across U.S.\ surgical programs due to variation in patient acuity, surgical practice, and rhythm telemetry ascertainment.\cite{filardo2017telemetry,perezgrovas2023jaha} Naive calendar holdouts allow nearly every test hospital to appear during training, masking transport losses across unseen hospitals and inflating apparent accuracy.\cite{steyerberg2016prediction,riley2019external} Second, while tabular foundation models (TFMs) and multi-layer automated machine learning (AutoML) frameworks have set new tabular benchmarks,\cite{hollmann2025tabpfn,erickson2025tabarena,erickson2020autogluon,jaeger2026tabpfn35,qu2026kumotabular,kong2026tabfm} national surgical evaluations rarely compare them side-by-side under both \emph{strictly unseen hospitals} (0\% overlap) and \emph{seen hospitals} (100\% overlap), nor do they quantify the inference latency and GPU memory burden of in-context foundation models at clinical deployment scale.

Here we conduct a leakage-audited national benchmark across all 368,132 isolated-CABG procedures from 1,091 U.S.\ centers in the Society of Thoracic Surgeons Adult Cardiac Surgery Database (STS ACSD) v2.81.\cite{sts2014v281,shahian2018sts} Using strictly 207 pre-incision predictors and identical 50,000-patient training cohorts across both unseen-hospital and seen-hospital splits, we benchmark five leading binary tabular AI models---Google TabFM v1.0.0, NVIDIA KumoTabular-Large v2, PriorLabs TabPFN-3.5, tuned CatBoost, and AutoGluon-Tabular 1.5 (\cli{best\_quality})---against NEVERHUMAN, a proprietary pre-incision clinical risk model, evaluating discrimination, calibration, clinical net benefit, subgroup equity, and wall-clock training and inference efficiency.
